
\documentclass[11pt]{article}
\usepackage{amsmath, amssymb}
\usepackage{geometry}
\usepackage{booktabs}
\usepackage{hyperref}
\geometry{margin=1in}

\title{Structural-Aware Handwriting Similarity Model}
\author{}
\date{}

\begin{document}
\maketitle

\section{Objective}
Design a high-speed, geometry-first model for online Chinese handwriting recognition and scoring, optimized for:
\begin{itemize}
\item Structural similarity rather than exact matching
\item Continuous accuracy percentages
\item Educational and learning-game feedback
\item Edge-device deployment
\end{itemize}

The system outputs both a character prediction and a similarity-based accuracy score reflecting how close a user’s drawing is to the canonical structure.

\section{Input Representation}

\subsection{Raw Input}
Each handwritten character is captured as a sequence of pen points:
\[
P = \{(x_i, y_i, t_i, p_i)\}, \quad i = 1 \dots N
\]

\subsection{Normalization}
\begin{enumerate}
\item Douglas--Peucker stroke simplification
\item Resampling to fixed length $L$ (e.g.\ $L=128$)
\item Spatial normalization to $[0,1]^2$
\end{enumerate}

\subsection{Per-Point Feature Vector}
Each point is encoded as:
\[
f_i = [x_i, y_i, \Delta x_i, \Delta y_i, \sin\theta_i, \cos\theta_i, \kappa_i]
\]

Final input tensor:
\[
X \in \mathbb{R}^{L \times D}
\]

\section{Model Architecture}

\subsection{Backbone}
\begin{itemize}
\item 1D Convolutional layers for local stroke primitives
\item Single-layer GRU for temporal modeling
\item Global average pooling
\end{itemize}

Output embedding:
\[
E \in \mathbb{R}^{d}, \quad d \approx 128
\]

\section{Dual-Head Output}

\subsection{Character Classification}
Softmax over $C$ characters:
\[
\hat{y}_{char} \in \mathbb{R}^{C}
\]

Loss:
\[
\mathcal{L}_{char} = \text{CrossEntropy}(\hat{y}_{char}, y_{char})
\]

\subsection{Structural Similarity Head}
Predicted structural embedding:
\[
\hat{y}_{struct} \in [0,1]^R
\]

Cosine similarity:
\[
S_{struct} = \frac{\hat{y}_{struct} \cdot y_{struct}}{\|\hat{y}_{struct}\|\|y_{struct}\|}
\]

Structural loss:
\[
\mathcal{L}_{struct} = 1 - S_{struct}
\]

\section{Training Objective}
\[
\mathcal{L}_{total} = \mathcal{L}_{char} + \lambda \mathcal{L}_{struct}, \quad \lambda \in [0.05,0.2]
\]

\section{Scoring for Learning Games}
Final score:
\[
Score = w_1 S_{struct} + w_2 S_{geom} + w_3 S_{prop}
\]

\begin{center}
\begin{tabular}{ll}
\toprule
Score & Feedback \\
\midrule
$\geq 90\%$ & Excellent \\
$75$--$90\%$ & Good \\
$60$--$75\%$ & Acceptable \\
$<60\%$ & Needs practice \\
\bottomrule
\end{tabular}
\end{center}

\section{Model Compression}
Iterative Magnitude Pruning removes 90--95\% of weights, followed by weight rewinding and retraining, producing a sparse winning-ticket network suitable for edge deployment.

\section{Deployment}
\begin{itemize}
\item Export to ONNX
\item INT8 quantization
\item Sub-millisecond inference on mobile CPUs
\item Fully offline operation
\end{itemize}

\section{Summary}
The proposed model treats handwriting as geometry rather than pixels, supports graded similarity scoring, enables radical-level feedback, and is efficient enough for real-time educational applications.

\end{document}
